% ===== プログラミングI 教材 PDF 用プリアンブル（XeLaTeX 専用） =====

% --- 日本語フォント（fontspec + XeTeX 内蔵の CJK 改行）---
% ※ この環境には xeCJK/ctex が無いため、XeTeX の locale 改行機能を使う。
\usepackage{fontspec}
\setmainfont{Noto Sans CJK JP}
\setsansfont{Noto Sans CJK JP}
\setmonofont{Noto Sans Mono CJK JP}[Scale=0.92] % ASCII も日本語コメントも表示可
% 日本語の行分割（単語間スペースの無い CJK を適切に折り返す）
\XeTeXlinebreaklocale "ja"
\XeTeXlinebreakskip = 0pt plus 0.1pt minus 0.01pt

% --- 余白 ---
\usepackage[a4paper,margin=18mm,headheight=14pt]{geometry}

% --- 見出し色 ---
\usepackage{xcolor}
\definecolor{headcol}{HTML}{1E3A5F}
\definecolor{subcol}{HTML}{2563A8}
\usepackage{sectsty}
\sectionfont{\color{headcol}\large}
\subsectionfont{\color{subcol}\normalsize}
\subsubsectionfont{\color{subcol}\normalsize}

% --- コードブロックの背景（pandoc の Shaded 環境）---
\definecolor{shadecolor}{HTML}{F1F5F9}
\setlength{\fboxsep}{4pt}

% --- 引用（blockquote）を左罫線つきに ---
\usepackage{framed}
\let\oldquote\quote
\let\endoldquote\endquote
\renewenvironment{quote}{%
  \begin{leftbar}\small\color{black!75}}{%
  \end{leftbar}}

% --- ヘッダ・フッタ ---
\usepackage{fancyhdr}
\pagestyle{fancy}
\fancyhf{}
\lhead{\small プログラミング I 演習資料}
\rhead{\small 2026 年度}
\cfoot{\small \thepage}
\renewcommand{\headrulewidth}{0.4pt}
\renewcommand{\footrulewidth}{0pt}

% --- 表の行間を少し詰める ---
\renewcommand{\arraystretch}{1.15}

% --- URL/長語の途中改行を許可 ---
\sloppy
